\scanpage{26}
sous l'hypothèse anabélienne faite, et celle de dominance. (Cette deuxième n'est pas superflue, comme on voit en prenant \deleted{pour $Y$ \inserted{$=G\backslash Z$, où $G$ est} un groupe elliptique et $Z$ un sous-groupe fini $\unsure{\ne 1|+1}$, et pour $f,g$} \inserted{deux applications constantes de $X$ dans $Y$!})

Tout d'abord, on vérifie qu'une variété élémentaire d'Artin se plonge \inserted{par un $i:Y\longrightarrow G$} bien dans une ext. d'une V.A. par un tore -- donc on aura $(i\circ g)=\tau_a(i\circ f)$. \deleted{Donc \unsure{tout}} \inserted{D'autre part,} \deleted{Quitte à remplacer} $f$,\deleted{$g$} étant dominante, il existe une partie ouverte \inserted{non vide} $\unsure{V|W}\subset Y$ contenue dans $\operatorname{Im}f$, et on aura donc $\unsure{V|W}+a\subset Y$. Soit $U\subset X$ un ouvert non vide tel que $f(U)\subset V$ \deleted{$g(U)\subset V$} \deleted{\unsure{donc qui} $\unsure{Q}$} quitte à remplacer $X$ par $U$, on est ramené au cas où $f$\deleted{,$g$} envoie $X$ dans un ouvert $V$ de $Y$ tel que $V+a\subset Y$, et $g\circ i=\tau_a(i\circ f)$. On \deleted{sait} \inserted{\unsure{voit}} que $\pi_1(f)$ est un composé (où $f':X\longrightarrow U$ est \deleted{\unsure{$f$}} induit par $f$ et $j:U\longrightarrow Y$ est l'inclusion)
\diagram{\begin{tikzcd}[ampersand replacement=\&]\pi_1(X) \arrow[r,"\pi_1(f')"] \& \pi_1(U) \arrow[r,"\pi_1(j)"] \& \pi_1(Y)\end{tikzcd}}
tandis que $\pi_1(g)$ est le composé
\diagram{\begin{tikzcd}[ampersand replacement=\&]\pi_1(X) \arrow[r,"\pi_1(f')"] \& \pi_1(U) \arrow[r,"{\pi_1(\unsure{j'|g'})}"] \& \pi_1(Y)\end{tikzcd}}
où $\deleted{\unsure{j}}$ $j':U\longrightarrow Y$ est induit par $\tau_a$.
\margin{25}
